\begin{table*}[!htb]
\centering
\caption{Data sources. Unity and Ski-PTZ-Pose support quantitative evaluation; the author-collected recordings are used only for qualitative illustration. Split counts are train / validation / test camera-pair sequences for Unity (action-disjoint) and camera-pair entries for Ski-PTZ-Pose.}
\label{tab:dataset_summary}
\footnotesize
\setlength{\tabcolsep}{4pt}
\renewcommand{\arraystretch}{1.12}
\begin{tabular}{@{}>{\raggedright\arraybackslash}p{0.135\linewidth}
                >{\raggedright\arraybackslash}p{0.30\linewidth}
                >{\raggedright\arraybackslash}p{0.26\linewidth}
                >{\raggedright\arraybackslash}p{0.235\linewidth}@{}}
\toprule
 & Unity synthetic & Ski-PTZ-Pose & Author-collected videos \\
\midrule
Type & Synthetic, rendered in Unity & Real, public giant-slalom footage & Real, indoor and outdoor recordings \\
Cameras & 180 synchronized views per action (5 layers \(\times\) 36 azimuths) & 6 synchronized PTZ cameras & 2 cameras per recording \\
Sequences & 24 (2 avatars \(\times\) 12 actions), 20--65 frames each & 12 runs, about 140 frames per camera (about 10K images) & 6 recordings (12 videos), 360--2,609 frames per view \\
Input & \multicolumn{3}{l}{Two monocular 3D pose streams from the native SAM3D front end} \\
Reference & Avatar skeleton, 15 joints & Canonical reference constructed from the annotations, 13 joints & None \\
Split & 193,320 / 128,880 / 64,440 & 150 / -- / 30 & Not used \\
Role & Training, main comparison, all controlled studies & Real-world evaluation & Qualitative deployment only \\
\bottomrule
\end{tabular}
\end{table*}
